% !TeX root = ./main.tex
% 项目排版覆盖：把 sysuthesis 输出对齐《中山大学航空航天学院学位论文撰写模板（2026年11月更新）》。
% 不改动 sysuthesis.cls；本文件由 main.tex 在 sysusetup.tex 之后载入，用于覆盖文档类默认值。
\makeatletter

% =====================================================================
% 一、章节标题
%   一级（第1章）  ：黑体小二(18bp)加粗、居中，段前段后各 1 行
%   二级（1.1）    ：黑体小三(15bp)加粗、左对齐，段前段后各 0.5 行
%   三级（1.1.1）  ：黑体四号(14bp)加粗、左对齐，段前段后各 0.5 行
%   四级（1.1.1.1）：宋体小四(12bp)加粗、左对齐，段前段后 0
% =====================================================================
\ctexset{
  chapter = {
    format      = \centering\heiti\bfseries\fontsize{18bp}{21.6bp}\selectfont,
    nameformat  = {},
    titleformat = {},
    number      = \thechapter,
    aftername   = \quad,
    beforeskip  = 24bp,
    afterskip   = 24bp,
    fixskip     = true,
  },
  section = {
    format     = \heiti\bfseries\fontsize{15bp}{21bp}\selectfont,
    aftername  = \quad,
    beforeskip = 12bp,
    afterskip  = 12bp,
  },
  subsection = {
    format     = \heiti\bfseries\fontsize{14bp}{21bp}\selectfont,
    aftername  = \quad,
    beforeskip = 12bp,
    afterskip  = 12bp,
  },
  subsubsection = {
    format     = \rmfamily\bfseries\fontsize{12bp}{20bp}\selectfont,
    number     = \thesubsection.\arabic{subsubsection},
    aftername  = \quad,
    indent     = 0em,
    beforeskip = 0bp,
    afterskip  = 0bp,
    runin      = false,
  },
}

% 目录：列出到三级标题（1.1.1），并在中文目录之后另附英文目录 Contents。
% 说明：\@starttoc 读取后会清空 .toc，故这里改用 \@input 读取上一轮生成的 .toc，
%       连续排版两遍，最后再打开写流供正文写入（等价于 \@starttoc 的写流部分）。
\renewcommand\tableofcontents{%
  \cleardoublepage
  \pagestyle{plain}%
  \sysu@chapter{\contentsname}%
  \setcounter{tocdepth}{2}%
  \begingroup\makeatletter\@input{\jobname.toc}\endgroup
  \cleardoublepage
  \pagestyle{plain}%
  \sysu@chapter{Contents}%
  \setcounter{tocdepth}{2}%
  \begingroup\makeatletter\@input{\jobname.toc}\endgroup
  \begingroup
    \makeatletter
    \if@filesw
      \expandafter\newwrite\csname tf@toc\endcsname
      \immediate\openout\csname tf@toc\endcsname\jobname.toc\relax
    \fi
    \@nobreakfalse
  \endgroup
}

% =====================================================================
% 二、正文行距：固定 20 磅（Word 模板“正文”样式为宋体小四、固定 20 磅）
% =====================================================================
\renewcommand\normalsize{%
  \@setfontsize\normalsize{12bp}{20bp}%
  \abovedisplayskip 6bp \@plus3bp \@minus3bp%
  \abovedisplayshortskip 6bp \@plus3bp \@minus3bp%
  \belowdisplayshortskip \abovedisplayshortskip
  \belowdisplayskip \abovedisplayskip
  \let\@listi\@listI}%
\normalsize
\ifx\MakeRobust\@undefined \else \MakeRobust\normalsize \fi

% =====================================================================
% 三、图表题注与目录名称
% =====================================================================
% 题注：五号宋体 / Times New Roman，不加粗（覆盖文档类中研究生图题的 font+=bf）。
\captionsetup{font=sysu}
% 英文图题标签用“Fig.”，表格仍用“Table”。
\renewcommand\sysu@set@names{%
  \ifsysu@language@chinese
    \def\figurename{图}%
    \def\tablename{表}%
  \else
    \def\figurename{Fig.}%
    \def\tablename{Table}%
  \fi
}
\sysu@set@names
% 图表目录名称：图目录 / 表目录。
\def\listfigurename{图目录}
\def\listtablename{表目录}

% =====================================================================
% 四、摘要与英文摘要
%   标题随章标题（小二）；正文固定行距 20 磅；移除抬头信息块；
%   关键词标签黑体加粗、词间以分号分隔；英文标签为“Keywords”。
% =====================================================================
\renewenvironment{abstract}{%
  \sysusetup{language=chinese}%
  \let\clearpage\relax
  \sysu@chapter{摘\hspace{2em}要}%
  \fontsize{12bp}{20bp}\selectfont
}{%
  \par\nobreak\smallskip\noindent\hangindent=4em\relax
  \textbf{\heiti 关键词}：\sysu@keywords@text\par
  \sysu@reset@main@language
  \newpage
}

\renewenvironment{abstract*}{%
  \sysusetup{language=english}%
  \let\clearpage\relax
  \ctexset{chapter/format=\centering\heiti\bfseries\fontsize{18bp}{21.6bp}\selectfont}%
  \sysu@chapter{ABSTRACT}%
  \fontsize{12bp}{20bp}\selectfont
}{%
  \par\nobreak\smallskip\noindent\hangindent=4em\relax
  \textbf{Keywords}: \sysu@keywords@en@text\par
  \newpage
  \sysu@reset@main@language
}

% =====================================================================
% 五、页眉名称（Word 模板为“中山大学工学博士学位论文”，此处按工学硕士填写）
% =====================================================================
\renewcommand\sysu@thesis@name{中山大学工学硕士学位论文}

% =====================================================================
% 七、封面与扉页
%   封面：密级/编号、校徽、学位论文类型、中英文题目、学位申请人/导师姓名及职称/专业名称、日期
%   扉页：中英文题目、专业/学位申请人/指导教师/学位、答辩委员会、日期
% =====================================================================

% 未设置密级时按公开处理。
\newcommand\sysu@secret@level@zh{%
  \ifdefined\sysu@secret@level
    \ifx\sysu@secret@level\@empty 公开\else\sysu@secret@level\fi
  \else
    公开
  \fi
}

% 封面正中的学位论文类型（Word 模板为“硕/博士学位论文”）。
\newcommand\sysu@cover@degree@name{%
  \ifsysu@degree@doctor 博士学位论文\else 硕士学位论文\fi
}

% 扉页日期格式：二〇二六年X月（汉字、无日）。
\newcommand\sysu@date@format@zh@cn[3]{\zhdigits{#1}年\zhnumber{#2}月}
\newcommand\sysu@date@zh@cn{\sysu@format@date{\sysu@date@format@zh@cn}{\sysu@date}}

% 研究生中文封面
\renewcommand\sysu@title@page@graduate@zh{%
  \sysusetup{language=chinese}%
  \begin{titlepage}%
    \sysu@pdfbookmark{封面}%
    \centering
    \parbox[t][0.6cm][t]{\textwidth}{%
      \raggedleft\heiti\fontsize{12bp}{12bp}\selectfont
      密级：\sysu@secret@level@zh\hspace{2em}编号：\sysu@student@id\par}\par
    \vskip 1.8cm%
    \includegraphics[width=8cm]{\sysu@logo@path/SYSULogo}\par
    \vskip 1.0cm%
    {\heiti\bfseries\fontsize{26bp}{30bp}\selectfont \sysu@cover@degree@name\par}
    \vskip 2.4cm%
    {\heiti\fontsize{22bp}{32bp}\selectfont \sysu@title\par}
    \vskip 0.8cm%
    {\rmfamily\fontsize{22bp}{30bp}\selectfont \sysu@title@en\par}
    \vskip 2.4cm%
    {\zihao{4}\sffamily
      \begin{tabular}{@{}l@{\hspace{0.6em}}l@{}}
        学位申请人：     & \sysu@author \\
        导师姓名及职称： & \sysu@supervisor@names \\
        专业名称：       & \sysu@speciality \\
      \end{tabular}\par}
    \vfill
    {\zihao{-4}\centering \sysu@date@zh@small\par}
  \end{titlepage}%
  \sysu@reset@main@language
}

% 研究生中文扉页（含答辩委员会表格）
\newcommand\sysu@title@page@graduate@zh@inner{%
  \sysusetup{language=chinese}%
  \begin{titlepage}%
    \sysu@pdfbookmark{扉页}%
    \centering
    \vspace*{1.2cm}%
    {\heiti\fontsize{18bp}{30bp}\selectfont \sysu@title\par}
    \vskip 0.8cm%
    {\rmfamily\fontsize{18bp}{28bp}\selectfont \sysu@title@en\par}
    \vskip 2.2cm%
    {\zihao{4}\sffamily
      \begin{tabular}{@{}l@{\hspace{0.6em}}l@{}}
        专\hspace{2\ccwd}业： & \sysu@speciality \\
        学位申请人：          & \sysu@author \\
        指导教师：            & \sysu@supervisor@names \\
        学\hspace{2\ccwd}位： & \thesisdegreename \\
      \end{tabular}\par}
    \vskip 2.2cm%
    {\zihao{4}
      \begin{tabular}{ccc}
        论文答辩委员会主席：& \raisebox{-1.5ex}{\underline{\makebox[6cm][c]{}}}& \makebox[3cm][c]{}\\
        &~\\
        论文答辩委员会成员：& \raisebox{-1.5ex}{\underline{\makebox[6cm][c]{}}}& \makebox[3cm][c]{}\\
        &~\\
        {}& \raisebox{-1.5ex}{\underline{\makebox[6cm][c]{}}}& \makebox[3cm][c]{}\\
        &~\\
        {}& \raisebox{-1.5ex}{\underline{\makebox[6cm][c]{}}}& \makebox[3cm][c]{}\\
        &~\\
        {}& \raisebox{-1.5ex}{\underline{\makebox[6cm][c]{}}}& \makebox[3cm][c]{}
      \end{tabular}\par}
    \vfill
    {\zihao{-5}\centering \sysu@date@zh@cn\par}
  \end{titlepage}%
  \sysu@reset@main@language
}

% 研究生封面输出“封面 + 扉页”两页。
% 直接重定义 \maketitle（与文档类一致，仅追加扉页）。
\renewcommand\maketitle{%
  \newgeometry{margin=2.54cm}%
  \pagenumbering{Alph}%
  \pagestyle{empty}%
  \sysu@title@pagetrue
  \ifsysu@degree@graduate
    \sysu@title@page@graduate@zh
    \sysu@title@page@graduate@zh@inner
  \else
    \sysu@cover
    \sysu@title@page@bachelor
  \fi
  \cleardoublepage
  \restoregeometry
  \pagestyle{plain}%
  \sysu@title@pagefalse
}

\makeatother
